% Passes 7171--7186: E8/D4/H27 association geometry and q=9 local exclusion.
\subsection{The $E_8/W(3,3)$ fibre scheme, selected $D_4$ geometry, and the $H_{27}$ matter cover}
\label{sec:pass7171-7186-e8-d4-h27-q9}

The six-to-one $E_8$-root fibration of Section~\ref{sec:pass7163-7170-e8-hexagonal-lift} extends to the complete five-class root association scheme.  Ordering the relations by normalized inner product $2,1,0,-1,-2$, the valencies are
\[
  \boxed{1,56,126,56,1},
\]
and the simultaneous eigenmatrix is
\[
P=\begin{pmatrix}
1&56&126&56&1\\
1&28&0&-28&-1\\
1&8&-18&8&1\\
1&-2&0&2&-1\\
1&-4&6&-4&1
\end{pmatrix},
\]
with multiplicities $1,8,35,112,84$.  The positive-inner-product adjacency has five distinct eigenspaces and therefore generates the full Bose--Mesner algebra.

\paragraph{The selected ninety $D_4$ subsystems.}
There are $3150$ $D_4$ root subsystems in $E_8$.  The ninety selected by the W33 center-quads form a rank-five association scheme with nontrivial valencies $1,32,32,24$ and eigenmatrix
\[
\begin{pmatrix}
1&1&32&32&24\\
1&-1&8&-8&0\\
1&1&-4&-4&6\\
1&1&2&2&-6\\
1&-1&-4&4&0
\end{pmatrix}.
\]
Its unique valency-one relation is $D_4\leftrightarrow D_4^\perp$.  The share-one-W33-point relation has exactly eighty maximal $9$-cliques: forty point stars and forty neighborhood stars, interchanged by orthogonal complementation.  Their quotient reconstructs W33, and consequently
\[
  \boxed{\operatorname{Aut}(\mathcal D_{90})\cong\operatorname{Aut}(W33)\times C_2},
  \qquad |\operatorname{Aut}(\mathcal D_{90})|=103680.
\]
All Krein parameters are nonnegative, but no primitive-idempotent ordering is $Q$-polynomial.

Each selected $D_4$ has exactly three mutually orthogonal four-frames on its twelve antipodal root lines; each frame uses one antipodal pair from each of the four constituent $C_6$ fibres.  Its positive-inner-product root graph has spectrum
\[
  \boxed{8^1\oplus4^4\oplus0^9\oplus(-2)^8\oplus(-4)^2}.
\]

\paragraph{Integral glue and the twenty-seven ten-$D_4$ spreads.}
The W33 relation of a selected $D_4$ pair determines its lattice behavior.  Share-one-fibre pairs span rank six.  The two disjoint nonorthogonal relations span rank eight with index one in $E_8$.  Orthogonal complement pairs span $D_4\oplus D_4$ with index four, with diagonal discriminant glue
\[
  \boxed{\{(0,0),(v,v),(s,s),(c,c)\}\cong(C_2)^2}.
\]
The twenty-seven ten-$D_4$ spreads contain five such orthogonal pairs each; any pair of $D_4$ subsystems drawn from different orthogonal pairs already generates the full $E_8$ lattice.

On the forty-five orthogonal-pair coordinates, the twenty-seven spread-incidence words generate
\[
  \boxed{[45,21,5]_2},
\]
and its dual is
\[
  \boxed{[45,24,6]_2}.
\]
The latter has exactly $120$ minimum words, objectwise the previously certified Steiner/$K_{3,3}$ circuits.  The all-ones word lies in the $21$-space; its even subcode is a twenty-dimensional full-group module.  An explicit intertwiner identifies that module with the previously certified binary tritangent selector $V_{20}$, including the outer PGSp generator.  Therefore
\[
  \boxed{C_{\rm spread}=\mathbf 1\oplus V_{20}^{\rm trit}},
  \qquad
  \boxed{C_{\rm spread}^{\perp}=V_{24}^{\rm Steiner}}.
\]

Every selected pair in the disjoint cross-count-four relation gives an explicit unimodular $E_8$ chart.  There are $1080$ such $GL_8(\mathbb Z)$ charts, partitioned by the twenty-seven spreads into forty each.  Their ordinary coordinate transitions $X_jX_i^{-1}$ form a flat groupoid: closed products telescope to the identity.  Thus any nontrivial triality-valued holonomy would require additional frame identifications and is not produced by basis changes alone.

\paragraph{$E_6\times A_2$, affine area, and the Heisenberg matter graph.}
Choose one six-root fibre.  It is an $A_2$ subsystem.  The W33 shell $1+12+27$ becomes
\[
  \boxed{6_{A_2}+72_{E_6}+81+81}.
\]
Either $81$-sector projects to twenty-seven $E_6$ minuscule weights with three $A_2$ charges per weight, recovering the Schlaefli graph $\operatorname{srg}(27,16,10,8)$.  The twenty-seven W33 nonneighbor fibres themselves are not the twenty-seven weights: they group into nine triples and form a regular $C_3$ voltage cover of $K_9$.

After switching a spanning star to voltage zero, its twelve zero-holonomy triangles are precisely the twelve affine lines of $AG(2,3)$.  Up to affine relabelling and global sign the voltage is
\[
  \boxed{\psi(u,v)=\det(u,v)\in\mathbb F_3},
\]
so triangle holonomy is the oriented affine area $\det(v-u,w-u)$.  Choosing the sign convention of the repo's canonical qutrit Heisenberg chart yields
\[
 (u,z)(v,w)=(u+v,z+w-\det(u,v)).
\]
The resulting twenty-seven-vertex E8 matter-fibre graph is exactly
\[
  \boxed{\operatorname{Cay}\!\left(H_{27},\{(u,0):u\ne0\}\right)}.
\]
It has intersection array
\[
  \boxed{\{8,6,1;1,3,8\}},
\]
spectrum $8^1\oplus2^{12}\oplus(-1)^8\oplus(-4)^6$, and full automorphism group
\[
  \boxed{H_{27}{:}GL_2(3)},\qquad |\operatorname{Aut}|=1296.
\]
The vertex stabilizer is transitive on all three distance spheres, so the graph is distance-transitive.

\paragraph{Two phase structures remain distinct.}
The $E_8$ fibre cocycle has only odd $\mathbb Z_{12}$ triangle holonomies $1,3,9,11$.  A two-qutrit Weyl commutator cochain has edge phase in $4\mathbb Z/12$ and hence triangle holonomy in $\{0,4,8\}$.  Since triangle holonomy is coboundary invariant, these two cocycles are not cohomologous.  No identification is made merely because both use twelfth roots of unity.

\paragraph{Local q=9 exclusion.}
In the canonical $(1,3,5)$ anchor graph, the known residual $47$-clique splits as forty-two invertible states plus five rank-one states.  Exact branch-and-bound excludes a target $48$ after deleting any $d\le8$ members of the forty-two-state invertible core and allowing arbitrary compatible additions from all remaining normalized states.  Thus a hypothetical target-$48$ residual clique in this anchor graph must delete at least nine core states.  This is a local basin theorem only:
\[
  \boxed{\text{no claim }\alpha(W(3,9))=51\text{ is made.}}
\]
